\section*{Bab \ref{ch:b1:functions} --- Fungsi Baku}

\begin{solution}{exo:b1:functions:1}
$\arcsin\frac{\sqrt 3}{2} = \frac\pi3$;
$\;\arccos\bigl(-\frac12\bigr) = \frac{2\pi}{3}$;
$\;\arctan(-1) = -\frac\pi4$.

$\arcsin\bigl(\sin\frac{5\pi}{6}\bigr)$: di sini $\sin\frac{5\pi}{6} =
\frac12$, dan sudut di $\intcc{-\frac\pi2}{\frac\pi2}$ yang sinusnya
$\frac12$ adalah $\frac\pi6$ (bukan $\frac{5\pi}{6}$).

$\arccos\bigl(\cos\frac{7\pi}{4}\bigr)$: di sini $\cos\frac{7\pi}{4} =
\frac{\sqrt 2}{2}$, dan sudut di $\intcc{0}{\pi}$ yang kosinusnya demikian
adalah $\frac\pi4$.
\end{solution}

\begin{solution}{exo:b1:functions:2}
$f(x) = \arcsin(2x - 1)$ menuntut $-1 \leq 2x - 1 \leq 1$, yakni $x
\in \intcc{0}{1}$. Pada $\intoo{0}{1}$, aturan rantai dan
\cref{prop:b1:functions:arcderiv} memberikan
\[
f'(x) = \frac{2}{\sqrt{1 - (2x-1)^2}} = \frac{2}{\sqrt{4x - 4x^2}}
= \frac{1}{\sqrt{x(1 - x)}} .
\]
$g(x) = \arctan\sqrt x$ terdefinisi untuk $x \geq 0$, dan untuk $x > 0$:
\[
g'(x) = \frac{1}{1 + x} \cdot \frac{1}{2\sqrt x}
= \frac{1}{2\sqrt x\,(1 + x)} .
\]
\end{solution}

\begin{solution}{exo:b1:functions:3}
$\cosh x + \sinh x = \frac{\eu^x + \eu^{-x}}{2} + \frac{\eu^x -
\eu^{-x}}{2} = \eu^x$. Jadi untuk $n \in \Z$:
\[
(\cosh x + \sinh x)^n = (\eu^x)^n = \eu^{nx}
= \cosh nx + \sinh nx .
\]
\end{solution}

\begin{solution}{exo:b1:functions:4}
Dengan $y = \arcsin x$: $\cos 2y = 1 - 2\sin^2 y = 1 - 2x^2$, jadi
$\cos(2\arcsin x) = 1 - 2x^2$.

Dengan $y = \arctan x$: $\sin 2y = 2 \sin y \cos y = 2 \tan y \cos^2 y =
\frac{2\tan y}{1 + \tan^2 y}$, jadi $\sin(2\arctan x) =
\dfrac{2x}{1 + x^2}$.
\end{solution}

\begin{solution}{exo:b1:functions:5}
Dari yang paling lambat ke yang paling cepat:
\[
(\ln x)^{1000} \;\ll\; x^{100} \;\ll\; x^{\ln x} \;\ll\;
\eu^{\sqrt x} \;\ll\; \eu^{x/100} .
\]
Alasannya: $(\ln x)^{1000}/x^{100} \to 0$ menurut
\cref{prop:b1:functions:powerrules} (logaritma kalah oleh pangkat).
Lalu $x^{100} = \eu^{100\ln x}$ dan $x^{\ln x} = \eu^{(\ln x)^2}$: karena
$(\ln x)^2 - 100\ln x \to +\infty$, yang kedua menang. Lalu $x^{\ln x} =
\eu^{(\ln x)^2} \ll \eu^{\sqrt x}$ karena $(\ln x)^2/\sqrt x \to 0$
(logaritma kalah oleh pangkat $x^{1/4}$, yang dikuadratkan). Akhirnya $\sqrt x - x/100
\to -\infty$, jadi $\eu^{\sqrt x} \ll \eu^{x/100}$.
\end{solution}

\begin{solution}{exo:b1:functions:6}
Daerah asalnya: $x \neq \pm 1$, yakni tiga selang. Pada masing-masingnya,
\[
f'(x) = \frac{\bigl(\frac{2x}{1-x^2}\bigr)'}{1 +
\frac{4x^2}{(1-x^2)^2}}
= \frac{\frac{2(1-x^2) + 4x^2}{(1-x^2)^2}}
{\frac{(1-x^2)^2 + 4x^2}{(1-x^2)^2}}
= \frac{2(1 + x^2)}{(1 + x^2)^2} = \frac{2}{1 + x^2}
= (2\arctan x)' .
\]
Jadi $f(x) - 2\arctan x$ konstan pada setiap selang. Nilainya: di $x =
0$ berlaku $f(0) = 0$, jadi konstantanya $0$ pada $\intoo{-1}{1}$. Ketika $x \to
+\infty$, $\frac{2x}{1-x^2} \to 0^-$ sehingga $f \to 0$, sedangkan $2\arctan x
\to \pi$: jadi konstantanya $-\pi$ pada $\intoo{1}{+\infty}$. Menurut keganjilannya,
konstanta itu $+\pi$ pada $\intoo{-\infty}{-1}$. Ringkasnya: $f = 2\arctan x$ pada
$\intoo{-1}{1}$, $f = 2\arctan x - \pi$ untuk $x > 1$, dan $f = 2\arctan x +
\pi$ untuk $x < -1$. (Inilah rumus sudut ganda bagi tangen,
yang dibaca lewat $\arctan$.)
\end{solution}

\begin{solution}{exo:b1:functions:7}
$\cosh x = 2$: dengan $u = \eu^x > 0$, $u + u^{-1} = 4$, jadi $u^2 - 4u +
1 = 0$, $u = 2 \pm \sqrt 3$: sehingga $x = \ln(2 + \sqrt 3)$ atau $x = \ln(2 -
\sqrt 3) = -\ln(2 + \sqrt 3)$ (kedua penyelesaiannya berlawanan, karena
$\cosh$ genap; keduanya sah). Setara dengan itu, $x = \pm
\operatorname{arcosh} 2$.

$5\cosh x - 4\sinh x = 3$: dengan mensubstitusikan definisi eksponensialnya,
$\frac{5(u + u^{-1}) - 4(u - u^{-1})}{2} = 3$, yakni $u + 9u^{-1} =
6$, yakni $u^2 - 6u + 9 = (u - 3)^2 = 0$: jadi $u = 3$, dengan satu penyelesaian $x =
\ln 3$.
\end{solution}

\begin{solution}{exo:b1:functions:8}
Misalkan $\alpha = \arctan\frac12 + \arctan\frac13$. Rumus penjumlahan:
\[
\tan\alpha = \frac{\frac12 + \frac13}{1 - \frac12\cdot\frac13}
= \frac{5/6}{5/6} = 1 .
\]
Kedua \omterm{def:b1:functions:arc}{arkus tangen} itu terletak di $\intoo{0}{\frac\pi4}$ (karena argumennya di
$\intoo{0}{1}$), jadi $\alpha \in \intoo{0}{\frac\pi2}$; dan satu-satunya sudut
di sana yang bertangen $1$ adalah $\frac\pi4$.

Untuk Machin: misalkan $\beta = \arctan\frac15$. Sudut ganda dua kali:
\[
\tan 2\beta = \frac{2/5}{1 - 1/25} = \frac{5}{12},
\qquad
\tan 4\beta = \frac{2 \cdot 5/12}{1 - 25/144} = \frac{120}{119}.
\]
Lalu
\[
\tan\Bigl(4\beta - \frac\pi4\Bigr)
= \frac{\frac{120}{119} - 1}{1 + \frac{120}{119}}
= \frac{1/119}{239/119} = \frac{1}{239}.
\]
Letaknya: $\beta < \arctan 1 = \frac\pi4$ dan sesungguhnya $\tan 4\beta =
\frac{120}{119}$ dekat dengan $1$ dengan $4\beta \in \intoo{0}{\frac\pi2}$
(karena $\beta < \frac\pi8$, sebab $\tan\frac\pi8 = \sqrt 2 - 1 > \frac15$);
jadi $4\beta - \frac\pi4 \in \intoo{-\frac\pi4}{\frac\pi4}$, yang di sana
$\arctan$ membalikkan $\tan$: $4\beta - \frac\pi4 =
\arctan\frac{1}{239}$, dan itulah rumus Machin.
\end{solution}

\begin{solution}{exo:b1:functions:9}
Misalkan $f(x) = x - \sin x$: $f(0) = 0$ dan $f'(x) = 1 - \cos x \geq 0$,
jadi $f \geq 0$ pada $\R_+$: sehingga $\sin x \leq x$.

Misalkan $g(x) = \sin x - x + \frac{x^3}{6}$: $g(0) = 0$, $g'(x) = \cos x -
1 + \frac{x^2}{2}$, $g'(0) = 0$, $g''(x) = -\sin x + x = f(x) \geq 0$
pada $\R_+$. Jadi $g'$ naik dengan $g'(0) = 0$, sehingga $g' \geq 0$,
sehingga $g$ naik dengan $g(0) = 0$: jadi $g \geq 0$, yakni $x -
\frac{x^3}{6} \leq \sin x$ pada $\R_+$.

Akibatnya $0 \leq \frac{x - \sin x}{x^3} \leq \frac 16$ untuk $x >
0$: rasionya terperangkap pada skala $\frac16$, dan
\cref{ch:b1:taylor} menunjukkan limitnya tepat $\frac 16$.
\end{solution}

\begin{solution}{exo:b1:functions:10}
Tulis $y = \arcsin x$, sehingga $2x\sqrt{1 - x^2} = 2\sin y\cos y = \sin
2y$ dan $h(x) = \arcsin(\sin 2y) - 2y$.

Untuk $\abs x \leq \frac{\sqrt 2}{2}$: $2y \in
\intcc{-\frac\pi2}{\frac\pi2}$, jadi $\arcsin(\sin 2y) = 2y$ dan $h =
0$.

Untuk $x \in \intoc{\frac{\sqrt 2}{2}}{1}$: $2y \in
\intoc{\frac\pi2}{\pi}$, dan sudut di
$\intcc{-\frac\pi2}{\frac\pi2}$ yang sinusnya $\sin 2y$ adalah $\pi - 2y$:
jadi $h(x) = \pi - 4\arcsin x$. (Periksa lewat turunannya: di sana
$h'(x) = \frac{2(1 - 2x^2)}{\abs{1 - 2x^2}\sqrt{1 - x^2}} -
\frac{2}{\sqrt{1-x^2}} = \frac{-4}{\sqrt{1 - x^2}}$, yaitu turunan
$-4\arcsin x$; dan di $x = 1$, $h(1) = \arcsin 0 - 2\cdot\frac\pi2
= -\pi = \pi - 4\cdot\frac\pi2$.)

Untuk $x \in \intco{-1}{-\frac{\sqrt 2}{2}}$, menurut keganjilan $h$:
$h(x) = -\pi - 4\arcsin x$.
\end{solution}

\begin{solution}{exo:b1:functions:11}
$g = \arctan \circ \sinh$ adalah komposisi fungsi ganjil yang naik
tegas, jadi ia ganjil dan naik tegas; ketika $x \to
+\infty$, $\sinh x \to +\infty$ sehingga $g(x) \to \frac\pi2$, dan $g$ adalah
bijeksi pada $\intoo{-\frac\pi2}{\frac\pi2}$ (dari kekontinuannya beserta
teorema nilai antara). Aturan rantai bersama
\cref{prop:b1:functions:hyprules}~(1):
\[
g'(x) = \frac{\cosh x}{1 + \sinh^2 x} = \frac{\cosh x}{\cosh^2 x}
= \frac{1}{\cosh x} .
\]
Menurut konstruksinya, $\tan g(x) = \sinh x$. Lalu, karena $g(x) \in
\intoo{-\frac\pi2}{\frac\pi2}$ berkosinus positif,
\begin{align*}
\cos g(x) &= \frac{1}{\sqrt{1 + \tan^2 g(x)}}
= \frac{1}{\sqrt{1 + \sinh^2 x}} = \frac{1}{\cosh x},\\
\sin g(x) &= \tan g(x) \cos g(x) = \frac{\sinh x}{\cosh x} = \tanh x .
\end{align*}
\end{solution}

\begin{solution}{exo:b1:functions:12}
Tulis $u = \arctan 2$ dan $v = \arctan 3$; keduanya terletak di
$\intoo{\frac\pi4}{\frac\pi2}$ (karena argumennya melebihi $1$), jadi $u
+ v \in \intoo{\frac\pi2}{\pi}$. Rumus penjumlahan memberikan
\[
\tan(u + v) = \frac{2 + 3}{1 - 6} = -1 ,
\]
dan satu-satunya sudut di $\intoo{\frac\pi2}{\pi}$ yang bertangen $-1$
adalah $\frac{3\pi}4$: jadi $\arctan 2 + \arctan 3 = \frac{3\pi}4$.
(Menerapkan $\arctan$ secara membabi buta pada tangennya akan memberikan
$-\frac\pi4$, meleset sejauh $\pi$ --- pelokalanlah yang menyelamatkan
perhitungan itu, dan penyebut $1 - 6$ yang negatif justru menjadi
isyarat bahwa jumlahnya sudah meninggalkan
$\intoo{-\frac\pi2}{\frac\pi2}$.) Dengan menambahkan $\arctan 1 = \frac\pi4$:
\[
\arctan 1 + \arctan 2 + \arctan 3 = \frac\pi4 + \frac{3\pi}4 = \pi .
\]
\end{solution}

\begin{solution}{pb:b1:functions:1}
\textbf{1.} Di sini $f$ adalah hasil bagi fungsi yang dapat diturunkan dengan
penyebut yang tak lenyap pada $\intoo0{+\infty}$, dan
\[
f'(t) = \frac{\frac1t \cdot t - \ln t}{t^2} = \frac{1 - \ln
t}{t^2} ,
\]
yang positif untuk $t < \eu$, nol di $t = \eu$, dan negatif untuk $t > \eu$:
jadi $f$ naik tegas pada $\intoc0\eu$, turun tegas pada
$\intco\eu{+\infty}$, dengan maksimum $f(\eu) = \frac1\eu$.

\textbf{2.} Ketika $t \to 0^+$: $\ln t \to -\infty$ dan $\frac1t \to
+\infty$, jadi $f(t) \to -\infty$. Ketika $t \to +\infty$: $f(t) \to 0$
menurut \omterm{prop:b1:functions:powerrules}{perbandingan pertumbuhan} pada \cref{prop:b1:functions:powerrules}
(dengan $\beta = \alpha = 1$). Tandanya: sama dengan tanda $\ln t$, jadi $f < 0$ pada
$\intoo01$, $f(1) = 0$, dan $f > 0$ pada $\intoo1{+\infty}$. Grafiknya
memanjat dari $-\infty$, memotong nol di $1$, memuncak di $(\eu,
\frac1\eu)$, lalu meluruh ke $0^+$.

\textbf{3.} $f(4) = \frac{\ln 4}4 = \frac{2\ln 2}4 = \frac{\ln 2}2
= f(2)$.

\textbf{4.} Menurut tabel variasinya dan teorema nilai
antara (yang dipakai pada taraf sekolah menengah; diformalkan pada
\cref{ch:b1:continuity}). Untuk $c < 0$: penyelesaian hanya ada di tempat
$f < 0$, yakni di $\intoo01$ yang di sana $f$ merupakan bijeksi naik
tegas pada $\intoo{-\infty}0$: jadi tepat satu. Untuk $c = 0$: hanya
$t = 1$. Untuk $0 < c < \frac1\eu$: pada $\intoo1\eu$, $f$ naik
dari $0$ ke $\frac1\eu$, jadi satu penyelesaian; pada $\intoo\eu{+\infty}$,
$f$ turun dari $\frac1\eu$ ke $0$, jadi satu lagi; totalnya dua. Untuk
$c = \frac1\eu$: hanya titik maksimumnya $t = \eu$. Untuk $c >
\frac1\eu$: tak ada.

\textbf{5.} Untuk $x \neq \eu$, kemaksimalan tegasnya memberikan $f(x) <
f(\eu)$, yakni $\frac{\ln x}x < \frac1\eu$, yakni $\eu\ln x < x$,
yakni $\ln(x^\eu) < \ln(\eu^x)$: jadi $x^\eu < \eu^x$. Dengan $x = \pi$:
$\eu^\pi > \pi^\eu$ (secara numerik $23.14 > 22.46$).

\textbf{6.} Untuk $x, y > 0$, kedua ruasnya positif, sehingga
\[
x^y = y^x \iff y\ln x = x\ln y \iff \frac{\ln x}x = \frac{\ln y}y
\iff f(x) = f(y) ,
\]
setelah dibagi $xy > 0$.

\textbf{7.} Misalkan $f(x) = f(y) = c$ dengan $x \neq y$. Jika $c \leq 0$ atau
$c = \frac1\eu$, pertanyaan 4 mengatakan persamaan $f = c$ hanya punya satu
penyelesaian: mustahil. Jadi $0 < c < \frac1\eu$, dan sekali lagi menurut
pertanyaan 4 kedua penyelesaiannya adalah satu titik di $\intoo1\eu$ dan satu titik
di $\intoo\eu{+\infty}$: jadi kedua koordinatnya melebihi $1$ dan keduanya
mengapit $\eu$.

\textbf{8.} Dengan mensubstitusikan $y = tx$ pada $y\ln x = x\ln y$:
\[
tx\ln x = x(\ln t + \ln x)
\iff (t - 1)\ln x = \ln t
\iff \ln x = \frac{\ln t}{t - 1} ,
\]
jadi $x = t^{1/(t-1)}$ dan $y = tx = t^{1 + \frac1{t-1}} =
t^{t/(t-1)}$. Sebaliknya, untuk nilai itu, $\ln y = \frac{t\ln
t}{t-1} = t\ln x$ dan $y = tx$, sehingga $\frac{\ln y}y = \frac{t\ln
x}{tx} = \frac{\ln x}x$: jadi sebuah penyelesaian. Setiap penyelesaian tak sepele punya
rasio $t = y/x \in \intoo0{+\infty} \setminus \{1\}$, jadi
parameterisasinya lengkap.

\textbf{9.} Untuk $t = 2$: $x = 2^{1/1} = 2$, $y = 2^{2/1} = 4$. Lalu
\[
x(1/t) = (1/t)^{\frac1{\frac1t - 1}}
= (t^{-1})^{\frac{t}{1 - t}}
= t^{\frac{t}{t - 1}} = y(t) ,
\]
lalu $y(1/t) = \frac1t\,x(1/t) = \frac{y(t)}t = x(t)$: jadi
pergantian parameter $t \mapsto 1/t$ menukarkan koordinatnya, sebagaimana
dituntut simetri persamaannya.

\textbf{10.} Ketika $t \to 1$: $\frac{\ln t}{t-1} \to 1$ (karena ia
hasil bagi selisih $\ln$ di $1$), jadi $x \to \eu^1 = \eu$ dan
$y = tx \to \eu$. Ketika $t \to +\infty$: $\ln x = \frac{\ln t}{t-1}
\to 0$ jadi $x \to 1$, sedangkan $\ln y = \frac{t\ln t}{t-1} \to
+\infty$ jadi $y \to +\infty$. Ketika $t \to 0^+$: $\ln x = \frac{\ln
t}{t-1} \to \frac{-\infty}{-1} = +\infty$ jadi $x \to +\infty$,
sedangkan $\ln y = \frac{t\ln t}{t-1} \to \frac{0}{-1} = 0$ jadi $y \to
1$ (dengan memakai $t\ln t \to 0$, \cref{ex:b1:functions:xx}). Jadi cabang
itu berjalan dari asimtot $y = 1$ (jauh di kanan), naik melewati
$(\eu, \eu)$ pada diagonalnya, lalu pergi menyusuri asimtot $x = 1$
(jauh di atas) --- dan simetris terhadap diagonalnya menurut pertanyaan 9.

\textbf{11.} $\ln x(t) = g(t) = \frac{\ln t}{t-1}$, dan
\[
g'(t) = \frac{\frac{t-1}t - \ln t}{(t-1)^2}
= \frac{h(t)}{(t-1)^2},
\qquad h(t) = 1 - \frac1t - \ln t .
\]
$h(1) = 0$ dan $h'(t) = \frac1{t^2} - \frac1t = \frac{1 -
t}{t^2} < 0$ untuk $t > 1$: jadi $h < 0$ pada $\intoo1{+\infty}$, sehingga $g' <
0$ dan $x = \eu^g$ turun tegas di sana (dari $\eu$ ke
$1$, menurut pertanyaan 10).

\textbf{12.} Menurut pertanyaan 11, $t \mapsto x(t)$ adalah bijeksi turun
tegas dari $\intoo1{+\infty}$ pada $\intoo1\eu$;
tulis $t(x)$ untuk inversnya (yang juga turun tegas) lalu tulis
$\varphi(x) = y(t(x))$. Perhatikan $h < 0$ juga pada $\intoo01$ (di sana
$h' > 0$ dan $h(1) = 0$), jadi $x(\cdot)$ turun pada
$\intoo01$ juga; sehingga $y(t) = x(1/t)$ (pertanyaan 9) bersifat
\emph{naik} dalam $t$ pada $\intoo1{+\infty}$, dari $\eu$ ke
$+\infty$. Dengan mengomposisikannya: $\varphi = y \circ t(\cdot)$ turun
tegas dari $\intoo1\eu$ pada $\intoo\eu{+\infty}$, dan
$x^{\varphi(x)} = \varphi(x)^x$ menurut pertanyaan 8. Akhirnya, untuk sebuah
nilai bersama $c = f(x) \in \intoo0{\frac1\eu}$, pasangan $\{x,
\varphi(x)\}$ adalah \emph{satu-satunya} \omterm{def:b1:logic:sets}{himpunan} penyelesaian $f = c$ yang beranggota dua
(pertanyaan 4); dengan memperluas $\varphi$ ke $\intoo\eu{+\infty}$ sebagai
\omterm{def:b1:logic:map}{pemetaan} inversnya, $\varphi(\varphi(x))$ kembali ke unsur yang lain, yakni
unsur asalnya: jadi $\varphi \circ \varphi = \mathrm{id}$.

\textbf{13.} Jika $(x, y)$ penyelesaian bulat yang tak sepele dengan $x
< y$, pertanyaan 7 menempatkan $x \in \intoo1\eu$: dan satu-satunya bilangan bulat
di sana adalah $x = 2$. Maka $f(y) = f(2) = \frac{\ln2}2$ dengan $y > \eu$; menurut
pertanyaan 4 persamaan $f = \frac{\ln2}2$ mempunyai tepat satu
penyelesaian di luar $\eu$, dan pertanyaan 3 menunjukkannya: $y = 4$. Jadi
$(2, 4)$, beserta tukarannya, adalah satu-satunya.

\textbf{14.} $t = 1 + \frac1n = \frac{n+1}n$ memberikan $\frac1{t - 1}
= n$ dan $\frac t{t-1} = n + 1$, jadi
\[
x_n = \Bigl(\frac{n+1}n\Bigr)^{n},
\qquad
y_n = \Bigl(\frac{n+1}n\Bigr)^{n+1},
\]
yang jelas rasional, berbeda (karena $y_n = t\,x_n \neq x_n$), dan $n =
1$ memberikan $x_1 = 2$, $y_1 = 4$.

\textbf{15.} Di sini $t = y/x$ rasional dan $> 1$; tulis $t - 1 =
\frac rs$ dalam bentuk paling sederhana, jadi $t = \frac{s + r}s$ dengan $\gcd(s +
r, s) = \gcd(r, s) = 1$. Pertanyaan 8 memberikan $x = t^{s/r}$, sehingga
\[
x^r = t^s = \frac{(s+r)^s}{s^s} ,
\]
yaitu pecahan dalam bentuk paling sederhana (tak ada prima yang membagi $s + r$ maupun
$s$). Dengan menulis $x = \frac pq$ dalam bentuk paling sederhana, $x^r = \frac{p^r}
{q^r}$ juga dalam bentuk paling sederhana, dan menurut ketunggalan penyajian
tersederhana itu: $p^r = (s+r)^s$ dan $q^r = s^s$. Bandingkan
eksponen sebarang prima $\ell$ pada $q^r = s^s$: $r \cdot
v_\ell(q) = s \cdot v_\ell(s)$, jadi $r$ membagi $s\,v_\ell(s)$;
karena $\gcd(r, s) = 1$, maka $r$ membagi $v_\ell(s)$ untuk setiap prima
$\ell$ (menurut ketunggalan pemfaktoran, yang diterima tanpa bukti; dibuktikan pada
\cref{ch:b1:arith}), sehingga $s = b^r$ untuk suatu bilangan bulat $b$. Argumen yang
sama pada $p^r = (s+r)^s$ memberikan $s + r = a^r$.

\textbf{16.} Andaikan $a^r - b^r = r$ dengan bilangan bulat $a > b \geq 1$
dan $r \geq 2$. Maka $a \geq b + 1$, jadi menurut teorema binomial
\[
a^r - b^r \geq (b+1)^r - b^r
= \sum_{k=0}^{r-1}\binom rk b^k
\geq \binom r0 + \binom r{r-1} b^{r-1}
\geq 1 + r > r ,
\]
yang merupakan kontradiksi. Dari pertanyaan 15, $a^r - b^r = (s + r) - s = r$
memaksa $r = 1$: jadi $t = 1 + \frac1s$, dan penyelesaiannya adalah $(x_s,
y_s)$ pada pertanyaan 14. Bersama pasangan yang tertukar,
penggolongannya pun lengkap.

\textbf{17.} $f\bigl(\tfrac94\bigr) = \frac{\ln 2.25}{2.25} =
\frac{0.81093}{2.25} = 0.3604$ dan $f\bigl(\tfrac{27}8\bigr) =
\frac{\ln 3.375}{3.375} = \frac{1.21640}{3.375} = 0.3604$ (sampai empat
angka desimal): sama, sebagaimana dijanjikan konstruksinya --- jadi
$\bigl(\tfrac94\bigr)^{27/8} = \bigl(\tfrac{27}8\bigr)^{9/4}$.

\textbf{18.} Parameternya $t_n = 1 + \frac1n$ turun tegas
ke $1$. Karena $x(\cdot)$ turun tegas pada
$\intoo1{+\infty}$ (pertanyaan 11), maka $x_n = x(t_n)$ \emph{naik}
tegas; dan karena $y(\cdot)$ naik tegas di sana
(pertanyaan 12), maka $y_n = y(t_n)$ \emph{turun} tegas. Menurut
pertanyaan 7, $x_n \in \intoo1\eu$ dan $y_n \in
\intoo\eu{+\infty}$: jadi $x_n < \eu < y_n$ untuk setiap $n$. Akhirnya
$t_n \to 1$ dan pertanyaan 10 memberikan $x_n \to \eu$ dan $y_n \to \eu$.
Jadi kekonvergenan monoton yang klasik dari $\bigl(1 +
\frac1n\bigr)^n$ jatuh begitu saja dari geometri cabang itu, tanpa
ketaksamaan baru.

\textbf{19.} Jika $\eu \leq a < b$: karena $f$ turun tegas pada
$\intco\eu{+\infty}$ maka $f(a) > f(b)$, yakni $b\ln a > a\ln
b$, yakni $a^b > b^a$. Jika $1 < a < b \leq \eu$: karena $f$ naik
tegas maka $f(a) < f(b)$, sehingga $a^b < b^a$. Kasus campurannya:
$2 < \eu < 3$ dan $2^3 = 8 < 9 = 3^2$; tetapi $2 < \eu < 5$ dan
$2^5 = 32 > 25 = 5^2$: jadi kedua hasilnya terjadi, dan yang memutuskan adalah
di sisi mana cabang itu titik $(a, b)$ jatuh.

\textbf{20.} Cabang itu bertemu diagonalnya di $\eu$ dan $\varphi$
diperluas secara kontinu di sana dengan $\varphi(\eu) = \eu$.
Menurunkan $\varphi(\varphi(x)) = x$ di $x = \eu$ dengan
aturan rantai: $\varphi'(\varphi(\eu))\,\varphi'(\eu) =
\varphi'(\eu)^2 = 1$, jadi $\varphi'(\eu) = \pm1$; dan karena $\varphi$
turun, maka $\varphi'(\eu) = -1$. Jadi cabang itu memotong
diagonalnya dengan kemiringan $-1$: yakni tegak lurus.

\textbf{21.} $y = 3x$ adalah kasus $t = 3$: $x = 3^{1/2} = \sqrt3$
dan $y = 3^{3/2} = 3\sqrt3$. Periksa: $f(\sqrt3) =
\frac{0.54931}{1.73205} = 0.3171$ dan $f(3\sqrt3) =
\frac{\ln 5.19615}{5.19615} = \frac{1.64792}{5.19615} = 0.3171$:
sama sampai empat angka desimal, jadi $(\sqrt3)^{3\sqrt3} =
(3\sqrt3)^{\sqrt3}$.

\textbf{22.} $n^{1/n} = \eu^{f(n)}$, dan $\exp$ naik, jadi
urutannya sama dengan urutan $f(n)$: $f(3) = \frac{\ln3}3 \approx
0.3662$ mengungguli $f(2) = f(4) = \frac{\ln2}2 \approx 0.3466$ (pertanyaan
3) dan $f(5) \approx 0.3219$. Jadi yang terbesar adalah $\sqrt[3]3$, dan
pasangan yang sama adalah $\sqrt2 = \sqrt[4]4$ --- yakni pasangan bulat $(2, 4)$
dalam samaran yang lain lagi.

\textbf{23.} \omterm{def:b1:logic:sets}{Himpunan} penyelesaiannya adalah gabungan diagonal $\{(x,
x) : x > 0\}$ dan satu cabang tunggal, yang simetris terhadap diagonalnya,
turun tegas dari asimtot $x = 1$ (ketika $y \to
+\infty$) ke asimtot $y = 1$ (ketika $x \to +\infty$), dan memotong
diagonalnya tepat sekali, di $(\eu, \eu)$, dengan kemiringan $-1$
di sana. Pada cabang itu duduk titik bulat $(2, 4)$ dan $(4, 2)$
--- satu-satunya --- beserta titik rasional $(x_n, y_n)$,
yang berbaris secara monoton di sepanjang cabang menuju $(\eu, \eu)$
tanpa pernah mencapainya (karena $\eu$ irasional; titik rasional cabang itu
menumpuk pada sebuah sudut yang irasional).

\textbf{24.} (i) \omterm{prop:b1:functions:powerrules}{Perbandingan pertumbuhan} memberikan limit $f$ di
$+\infty$ (pertanyaan 2) dan $t\ln t \to 0$ (pertanyaan 10),
yang membentuk tabel variasinya sekaligus asimtotnya. (ii) Teorema
nilai antara, lewat \omterm{def:b1:logic:statement}{pernyataan} bijeksinya,
mengubah tabel variasi itu menjadi cacahan penyelesaian yang eksak
(pertanyaan 4) dan menjadi keberadaan \omterm{def:b1:logic:map}{pemetaan} invers $t(x)$
(pertanyaan 12). (iii) Ketunggalan pemfaktoran menggerakkan kedua
penyamaan bentuk tersederhana pada pertanyaan 15, yang menjadi jantung aritmetika
bagi penggolongan titik rasionalnya.

\textbf{25.} Satu perhitungan turunan --- yaitu tanda $1 - \ln
t$ --- melahirkan segalanya: cacahan penyelesaian untuk setiap aras
$c$, bentuk dan asimtot cabangnya, ketaksamaan
ekstremal $\eu^x \geq x^\eu$, penggolongan penyelesaian bulat dan
rasional, serta kekonvergenan monoton $\bigl(1 +
\frac1n\bigr)^n$. Inilah penghematan berpikir dengan tabel
variasi: telaah berdimensi satu menyelesaikan persamaan dengan dua variabel
karena persamaannya melewati satu fungsi tunggal.
Nanti \cref{ch:b1:seq} akan mengulang kekonvergenan $(x_n)$ dengan
teori umum barisan monoton; \cref{ch:b1:derivative}
melandasi tabel variasi secara ketat di atas teorema nilai rata-rata; dan
\cref{ch:b1:taylor} memasok apa yang tak dapat diberikan tabel itu --- yaitu
\emph{laju} kekonvergenan $x_n \to \eu$ (yang berorde
$1/n$).

\end{solution}
